% ============================================================
% Arbeitsblatt für die Gruppen
% ============================================================
\clearpage
\setstretch{1.1}
\abschnitt{Arbeitsauftrag: Die Zukunft von Endstation 7}
\noindent Gruppe: \rule{5.5cm}{0.4pt}\hfill Delegation: \rule{5cm}{0.4pt}

\medskip
\noindent\textbf{Ihre Rolle.} Sie sind die anerkannten Vertreterinnen und Vertreter einer der fünf Gemeinschaften von Endstation~7. Ihre Gruppe kontrolliert einen Bereich, der für das Überleben aller Menschen notwendig ist. Die bisherige Stationsleitung ist zusammengebrochen. Sie müssen gemeinsam mit den anderen Gruppen eine neue politische Ordnung entwickeln.

\subsection*{Aufgabe 1: Bestandsaufnahme \hfill \normalsize\ifkurz 5\else 8\fi{} Min.}
\begin{enumerate}
  \item Analysieren Sie Ihre Situation: Welche Bedeutung haben Ihre Marker für das Überleben der Station?
  \item Welche Probleme kann Ihre Gemeinschaft ohne Zusammenarbeit nicht lösen?
  \item Legen Sie \textbf{drei Interessen} fest, die Sie unbedingt durchsetzen wollen.
  \item Welche \textbf{Zugeständnisse} könnten Sie anderen Gruppen machen?
\end{enumerate}
\linien{3}

\subsection*{Aufgabe 2: Entscheidung über die Regierungsform}
Welche Ordnung bevorzugen Sie für Endstation~7? Sie dürfen eine bekannte Regierungsform wählen, mischen oder ein eigenes Modell entwickeln (Infokarten).

\noindent
\feld Direkte Demokratie \quad \feld Repräsentative Demokratie \quad \feld Autoritäre Führung\\[1mm]
\feld Technokratie \quad \feld Föderation \quad \feld Mischform oder eigenes Modell: \rule{3.8cm}{0.4pt}

\medskip
\noindent Begründen Sie Ihre Wahl anhand von \textbf{mindestens drei Kriterien}: Wie schnell fallen Entscheidungen? Wie gerecht werden Interessen berücksichtigt? Wie wird verhindert, dass einzelne Gruppen zu viel Macht erhalten? Wie wird die Versorgung langfristig gesichert? Wie werden Regeln durchgesetzt und Konflikte gelöst?
\linien{3}

\ifkurz\else\clearpage\fi
\subsection*{Aufgabe 3: Verhandlungsvorbereitung \ifkurz\else\hfill \normalsize 5 Min.\fi}
Entwickeln Sie ein Angebot an die anderen Gruppen.
\par\smallskip
{\renewcommand{\arraystretch}{1}
\begin{tabular}{|p{4.6cm}|p{\dimexpr\textwidth-4.6cm-4\tabcolsep-3\arrayrulewidth\relax}|}
\hline
\rule{0pt}{\ifkurz 9mm\else 11mm\fi}Was können wir anbieten? & \\ \hline
\rule{0pt}{\ifkurz 9mm\else 11mm\fi}Was brauchen wir von anderen? & \\ \hline
\rule{0pt}{\ifkurz 9mm\else 11mm\fi}Welche Bedingungen stellen wir? & \\ \hline
\rule{0pt}{\ifkurz 9mm\else 11mm\fi}\textbf{Unsere rote Linie:} & \\ \hline
\rule{0pt}{\ifkurz 9mm\else 11mm\fi}Mit wem könnten wir eine Koalition bilden? & \\ \hline
\end{tabular}}

\subsection*{Aufgabe 4: Gemeinsamer Regierungsvertrag}
Verhandeln Sie mit den anderen Gruppen über eine gemeinsame Ordnung und halten Sie sie im \textbf{Regierungsvertrag von Endstation~7} fest. Ihr Ziel ist nicht, alle anderen Gruppen zu besiegen: Die Vereinbarung muss das Überleben der gesamten Station ermöglichen. Ihre Gruppe muss am Ende erklären können, warum sie zustimmt oder ablehnt.

\subsection*{Aufgabe 5: Rückblick \hfill \normalsize\ifkurz Hausaufgabe\else Hausaufgabe oder Sicherung\fi}
\begin{enumerate}
  \item Welche Interessen haben Ihre Entscheidung beeinflusst?
  \item Welche Vorteile und Nachteile hat die beschlossene Regierungsform?
  \item Wer konnte seine Interessen besonders gut durchsetzen?
  \item Ist die Vereinbarung gerecht, stabil und demokratisch?
  \item Was würden Sie bei einer erneuten Verhandlung anders machen?
\end{enumerate}
\linien{3}
